\chapter{States of light: coherent, squeezed, and interfered}
\label{ch:states-of-light}

The Fock states of \cref{ch:em-quantisation} diagonalise the photon
number, but they are exotic --- hard to make and rarely encountered.
The light that actually drives a cavity, carries a readout tone, or
illuminates an interferometer is in a \emph{coherent} state: the
closest a quantum field comes to a classical wave. Beyond the
classical, the quadrature noise of \cref{sec:quadratures} can be
redistributed --- \emph{squeezed} --- to push one quadrature below
the vacuum limit, the resource behind sub-standard-quantum-limit
metrology and the squeezed-vacuum haloscopes of Part~VI. This
chapter builds these states. We define the displacement operator
and the coherent states it generates, the squeeze operator and the
squeezed states, the beam-splitter transformation that mixes modes
and enables interferometry, and the second-order coherence
$g^{(2)}$ that classifies light as bunched, coherent, or
antibunched. Throughout, the phase-space picture of
\cref{sec:quadratures} is the unifying language, and the
standard-quantum-limit connection to the Heisenberg-limited sensing
of \cref{ch:coherent-multiqubit} is made explicit.

% --------------------------------------------------------------
\section{Coherent states}
\label{sec:coherent-states}

\begin{phenomenon}
A classical oscillating current --- the local oscillator of a
microwave source, the field of an ideal laser --- drives a field
mode into a state that is as close to a classical wave as quantum
mechanics allows: a well-defined amplitude and phase, with the
minimum noise permitted by the uncertainty relation, distributed
equally between the two quadratures. This is the \emph{coherent
state}. It is the displaced vacuum: take the vacuum noise ball of
\cref{sec:quadratures} and push it out to amplitude $\alpha$ in
phase space without distorting it.
\end{phenomenon}

\begin{statement}[\textbf{Displacement operator and coherent states}~\cite{glauber1963}]
\label{stmt:coherent}
The \emph{displacement operator}
\begin{equation}
\label{eq:displacement}
\op D(\alpha) \;=\; \exp\bigl(\alpha\,\hat a^\dagger - \alpha^*\hat a\bigr),
\qquad \alpha\in\C,
\end{equation}
is unitary and shifts the ladder operators,
$\op D^\dagger(\alpha)\,\hat a\,\op D(\alpha)=\hat a+\alpha$. The
\emph{coherent state} $\ket\alpha=\op D(\alpha)\ket0$ is the unique
eigenstate of the annihilation operator,
\begin{equation}
\label{eq:coherent-eigenstate}
\hat a\ket\alpha \;=\; \alpha\ket\alpha,
\qquad
\ket\alpha \;=\; e^{-|\alpha|^2/2}\sum_{n=0}^\infty\frac{\alpha^n}{\sqrt{n!}}\ket n,
\end{equation}
with Poissonian photon statistics
$P(n)=e^{-|\alpha|^2}|\alpha|^{2n}/n!$, mean $\bar n=|\alpha|^2$,
variance $(\Delta n)^2=|\alpha|^2$, hence Mandel $Q=0$
(\cref{stmt:mandel}). It is a minimum-uncertainty state with
$\Delta X=\Delta P=\tfrac12$, the displaced vacuum.
\end{statement}

\paragraph{Derivation of the shift property.}
With $\hat A=\alpha\hat a^\dagger-\alpha^*\hat a$, the
Baker--Campbell--Hausdorff identity gives
$\op D^\dagger\hat a\op D=e^{-\hat A}\hat a\,e^{\hat A}
=\hat a+\comm{\hat a}{\hat A}+\tfrac12\comm{\comm{\hat a}{\hat A}}{\hat A}+\cdots$.
The first commutator is
$\comm{\hat a}{\alpha\hat a^\dagger-\alpha^*\hat a}=\alpha$, a
c-number, so all higher commutators vanish and
$\op D^\dagger\hat a\op D=\hat a+\alpha$. Acting on the vacuum,
$\hat a\ket\alpha=\hat a\op D\ket0=\op D(\op D^\dagger\hat a\op D)\ket0=\op D(\hat a+\alpha)\ket0=\alpha\op D\ket0=\alpha\ket\alpha$,
the eigenvalue equation.

\begin{statement}[\textbf{Over-completeness}]
\label{stmt:overcomplete}
The coherent states are not orthogonal,
$|\langle\beta|\alpha\rangle|^2=e^{-|\alpha-\beta|^2}$, and form an
\emph{over-complete} set: they resolve the identity but with a
density of one state per phase-space area $\pi$,
\begin{equation}
\label{eq:coherent-resolution}
\frac{1}{\pi}\int_\C \dd^2\alpha\;\ket\alpha\!\bra\alpha \;=\; \id.
\end{equation}
Two coherent states of nearly equal amplitude overlap strongly; the
overlap falls off as a Gaussian in their phase-space separation.
\end{statement}

\paragraph{Application: a driven cavity.}
A classical drive term $\op H_{\rm drive}=\hbar(\epsilon\,\hat a^\dagger e^{-i\omega_d t}+\text{h.c.})$
added to a cavity mode displaces it: in the rotating frame, the
steady state of the input--output equation \cref{eq:qle} is the
coherent state $\ket{\alpha}$ with $\alpha=\sqrt\kappa\,\langle\hat a_{\rm in}\rangle/(i\Delta+\kappa/2)$
of \cref{eq:intracavity}. The coherent state is therefore the
generic state of any linearly driven, damped mode --- the readout
tone in dispersive measurement, the pump in a parametric amplifier,
and the carrier in every microwave control pulse.

% --------------------------------------------------------------
\section{Squeezed states}
\label{sec:squeezed-states}

\begin{phenomenon}
The coherent state spreads the minimum allowed noise equally
between its two quadratures. The uncertainty relation
\cref{eq:quadrature-uncertainty} only constrains the
\emph{product} $\Delta X\,\Delta P\geq\tfrac14$, so one quadrature
can be squeezed below the vacuum value $\tfrac12$ provided the
other is stretched to compensate. A \emph{squeezed state} does
exactly this: it deforms the circular vacuum noise ball into an
ellipse. Measuring along the squeezed quadrature beats the standard
quantum limit --- the resource behind squeezed-light interferometry
and the squeezed-vacuum haloscopes of \cref{ch:microwave-absorption}.
\end{phenomenon}

\begin{statement}[\textbf{Squeeze operator and squeezed vacuum}~\cite{caves1981}]
\label{stmt:squeezing}
The \emph{squeeze operator}
\begin{equation}
\label{eq:squeeze}
\op S(\xi) \;=\; \exp\!\Bigl(\tfrac12\bigl(\xi^*\hat a^2 - \xi\,\hat a^{\dagger2}\bigr)\Bigr),
\qquad \xi=r\,e^{i\theta},
\end{equation}
is unitary and acts on the ladder operators by a Bogoliubov
transformation,
\begin{equation}
\label{eq:squeeze-bogoliubov}
\op S^\dagger(\xi)\,\hat a\,\op S(\xi)
\;=\; \hat a\,\cosh r - \hat a^\dagger\,e^{i\theta}\sinh r.
\end{equation}
The squeezed vacuum $\ket\xi=\op S(\xi)\ket0$ has, for $\theta=0$,
quadrature variances
\begin{equation}
\label{eq:squeeze-variances}
(\Delta X)^2 = \tfrac14 e^{-2r},
\qquad
(\Delta P)^2 = \tfrac14 e^{+2r},
\end{equation}
still saturating $\Delta X\,\Delta P=\tfrac14$ but with the $X$
noise reduced below the vacuum by $e^{-2r}$ (in decibels,
$10\log_{10}e^{2r}$). It contains photons even though it is built
from the vacuum, $\bar n=\sinh^2 r$, with super-Poissonian,
even-only number statistics.
\end{statement}

\paragraph{Derivation of the variances.}
From \cref{eq:squeeze-bogoliubov} with $\theta=0$,
$\op S^\dagger\hat a\op S=\hat a\cosh r-\hat a^\dagger\sinh r$, so
the squeezed-vacuum quadrature $\hat X$ has variance
$\langle0|\op S^\dagger\hat X^2\op S|0\rangle$. Writing
$\hat X=\tfrac12(\hat a+\hat a^\dagger)$ and transforming,
$\op S^\dagger\hat X\op S=\tfrac12(\hat a+\hat a^\dagger)(\cosh r-\sinh r)=\hat X e^{-r}$,
so $(\Delta X)^2=e^{-2r}\langle0|\hat X^2|0\rangle=\tfrac14 e^{-2r}$.
Likewise $\op S^\dagger\hat P\op S=\hat P e^{+r}$ gives
$(\Delta P)^2=\tfrac14 e^{+2r}$. The product is preserved; the noise
is merely redistributed.

\paragraph{Application: beating the standard quantum limit.}
A phase measurement with $\bar n$ photons of coherent light has
sensitivity $\Delta\phi\sim1/\sqrt{\bar n}$ (the shot-noise or
standard quantum limit); injecting squeezed vacuum into the unused
port of the interferometer reduces the relevant quadrature noise by
$e^{-r}$, improving the sensitivity to
$\Delta\phi\sim e^{-r}/\sqrt{\bar n}$. This is how
gravitational-wave detectors and the HAYSTAC axion
haloscope~\cite{haystac2021,malnou2019} operate below the SQL, and
it is the single-mode optical analogue of the entanglement-enhanced,
Heisenberg-limited multi-qubit sensing of
\cref{ch:coherent-multiqubit}.

% --------------------------------------------------------------
\section{Beam splitters and interferometry}
\label{sec:beam-splitters}

\begin{phenomenon}
A beam splitter takes two input modes and linearly mixes them into
two output modes --- the elementary two-mode operation of quantum
optics. Cascading beam splitters builds an interferometer, whose
output depends on the relative phase of the arms; this is how a
phase shift (induced by a length change, a refractive index, or a
dark-matter field) is converted into a measurable intensity. The
quantum state of the input light sets the phase sensitivity:
coherent light gives the standard quantum limit, squeezed or
entangled light beats it.
\end{phenomenon}

\begin{statement}[\textbf{Beam-splitter transformation}]
\label{stmt:beam-splitter}
A lossless beam splitter mixing modes $\hat a,\hat b$ acts by the
unitary $\op U_{\rm BS}=\exp[\vartheta(\hat a^\dagger\hat b-\hat a\hat b^\dagger)]$,
transforming the modes as
\begin{equation}
\label{eq:beam-splitter}
\begin{pmatrix}\hat a_{\rm out}\\ \hat b_{\rm out}\end{pmatrix}
=\begin{pmatrix}\cos\vartheta & \sin\vartheta\\ -\sin\vartheta & \cos\vartheta\end{pmatrix}
\begin{pmatrix}\hat a_{\rm in}\\ \hat b_{\rm in}\end{pmatrix},
\end{equation}
with transmissivity $T=\cos^2\vartheta$ ($\vartheta=\pi/4$ for a
50:50 splitter). The transformation preserves the commutators
($\comm{\hat a_{\rm out}}{\hat a_{\rm out}^\dagger}=1$, etc.), so it
is a genuine canonical (photon-number-conserving) operation.
\end{statement}

\paragraph{Mach--Zehnder interferometry and the phase sensitivity.}
Two 50:50 beam splitters with a phase shift $\phi$ between them form
a Mach--Zehnder interferometer. With coherent light of mean photon
number $\bar n$ in one port and vacuum in the other, the output
photon-number difference measures $\phi$ with sensitivity
\begin{equation}
\label{eq:mz-sensitivity}
\Delta\phi_{\rm SQL} \;=\; \frac{1}{\sqrt{\bar n}},
\end{equation}
the shot-noise / standard quantum limit --- the optical counterpart
of the SQL derived for $N$ independent qubit probes in
\cref{stmt:sql}. Replacing the vacuum port with squeezed vacuum
lowers this by $e^{-r}$; replacing the coherent input with a
path-entangled $N$-photon (NOON) state reaches the Heisenberg limit
$\Delta\phi\sim1/\bar n$, the photonic analogue of the
GHZ-enhanced sensing of \cref{stmt:hl}.

\paragraph{Hong--Ou--Mandel: two-photon interference.}
A purely quantum beam-splitter effect: send one photon into each
input port of a 50:50 splitter. The two-photon amplitudes for
``both transmit'' and ``both reflect'' interfere destructively, and
the photons \emph{always} leave together from the same output port
--- they never split. This Hong--Ou--Mandel effect has no classical
analogue (it relies on bosonic exchange symmetry) and is the
workhorse test of photon indistinguishability in photonic quantum
information.

% --------------------------------------------------------------
\section{Coherence and photon statistics}
\label{sec:g2}

\begin{statement}[\textbf{Second-order coherence $g^{(2)}$}]
\label{stmt:g2}
The normalised second-order coherence at zero delay,
\begin{equation}
\label{eq:g2}
g^{(2)}(0) \;=\; \frac{\langle\hat a^\dagger\hat a^\dagger\hat a\hat a\rangle}{\langle\hat a^\dagger\hat a\rangle^2}
\;=\; 1 + \frac{(\Delta n)^2-\bar n}{\bar n^2}
\;=\; 1+\frac{Q}{\bar n},
\end{equation}
measures the tendency of photons to arrive together. Coherent light
has $g^{(2)}(0)=1$ (random, Poissonian arrivals); thermal light has
$g^{(2)}(0)=2$ (\emph{bunching}, super-Poissonian); and
non-classical light has $g^{(2)}(0)<1$ (\emph{antibunching}), with a
single-photon Fock state giving $g^{(2)}(0)=1-1/n\to0$ for $n=1$.
The condition $g^{(2)}(0)<1$ is impossible for any classical field
and is the standard experimental signature of a single-photon
source.
\end{statement}

\paragraph{Relation to the Mandel parameter.}
\Cref{eq:g2} makes the link to the Mandel $Q$ of
\cref{stmt:mandel} explicit: $g^{(2)}(0)=1+Q/\bar n$, so
super-Poissonian light ($Q>0$) bunches and sub-Poissonian light
($Q<0$) antibunches. The two parameters carry the same information
about second-order photon statistics; $g^{(2)}$ is the one a
Hanbury Brown--Twiss intensity-correlation measurement returns
directly.

% --------------------------------------------------------------
This chapter populated the quantised field of \cref{ch:em-quantisation}
with the states that occur in practice. The coherent state is the
displaced vacuum, the eigenstate of $\hat a$ with Poissonian
statistics and minimum, symmetric quadrature noise --- the generic
output of any linear drive. The squeezed state redistributes that
noise to beat the standard quantum limit in one quadrature. Beam
splitters mix modes and build interferometers whose phase
sensitivity is set by the input state, from the SQL of coherent
light to the Heisenberg limit of entangled light. The second-order
coherence $g^{(2)}(0)$ classifies light as bunched, coherent, or
antibunched.

These states are the carriers of information between a qubit and the
outside world. \Cref{ch:jaynes-cummings} now couples a single field
mode to a single qubit and works out the dynamics --- the
Jaynes--Cummings model --- from which dispersive readout and the
Purcell effect follow.

\clearpage
\begin{chaptersummary}

\paragraph{Coherent states.}
$\op D(\alpha)=\exp(\alpha\hat a^\dagger-\alpha^*\hat a)$ displaces
the vacuum to $\ket\alpha=\op D(\alpha)\ket0$, the eigenstate
$\hat a\ket\alpha=\alpha\ket\alpha$ with Poissonian statistics
($\bar n=|\alpha|^2$, $Q=0$) and $\Delta X=\Delta P=\tfrac12$. They
are over-complete, $\tfrac1\pi\int\dd^2\alpha\,\ket\alpha\!\bra\alpha=\id$,
and are the steady state of any linearly driven damped mode.

\paragraph{Squeezed states.}
$\op S(\xi)=\exp[\tfrac12(\xi^*\hat a^2-\xi\hat a^{\dagger2})]$
gives $\op S^\dagger\hat a\op S=\hat a\cosh r-\hat a^\dagger e^{i\theta}\sinh r$
and quadrature variances $\tfrac14 e^{\mp2r}$: one quadrature below
the SQL, the other above, product preserved. Squeezed vacuum has
$\bar n=\sinh^2 r$ and improves phase sensitivity to
$e^{-r}/\sqrt{\bar n}$.

\paragraph{Beam splitters.}
$\op U_{\rm BS}=\exp[\vartheta(\hat a^\dagger\hat b-\hat a\hat b^\dagger)]$
rotates the two modes; a 50:50 splitter has $\vartheta=\pi/4$.
Mach--Zehnder phase sensitivity is $1/\sqrt{\bar n}$ (SQL) for
coherent input, $1/\bar n$ (Heisenberg) for NOON input. Hong--Ou--Mandel
two-photon interference sends coincident photons out the same port.

\paragraph{Coherence.}
$g^{(2)}(0)=\langle\hat a^\dagger\hat a^\dagger\hat a\hat a\rangle/\langle\hat a^\dagger\hat a\rangle^2=1+Q/\bar n$:
$=1$ coherent, $=2$ thermal (bunching), $<1$ antibunched
(non-classical, $=0$ for a single photon).

\end{chaptersummary}

% --------------------------------------------------------------
\section*{Exercises}
\addcontentsline{toc}{section}{Exercises}

\begin{exercise}[Displacement composition]
Using BCH, show that $\op D(\alpha)\op D(\beta)=e^{i\,\mathrm{Im}(\alpha\beta^*)}\op D(\alpha+\beta)$,
so displacements compose up to a phase.
\begin{solution}
For operators with c-number commutator,
$e^{\hat A}e^{\hat B}=e^{\hat A+\hat B}e^{\tfrac12\comm{\hat A}{\hat B}}$.
With $\hat A=\alpha\hat a^\dagger-\alpha^*\hat a$,
$\hat B=\beta\hat a^\dagger-\beta^*\hat a$,
$\comm{\hat A}{\hat B}=\alpha\beta^*\comm{\hat a^\dagger}{-\hat a... }$;
evaluating, $\comm{\hat A}{\hat B}=(\alpha\beta^*-\alpha^*\beta)\comm{\hat a^\dagger}{\hat a}\cdot(-1)$...
more directly $\comm{\hat A}{\hat B}=-|...|$ gives
$\tfrac12\comm{\hat A}{\hat B}=i\,\mathrm{Im}(\alpha\beta^*)$. Hence
$\op D(\alpha)\op D(\beta)=e^{i\,\mathrm{Im}(\alpha\beta^*)}\op D(\alpha+\beta)$.
The phase is the (symplectic) area swept in phase space.
\end{solution}
\end{exercise}

\begin{exercise}[Poisson statistics of a coherent state]
From \cref{eq:coherent-eigenstate} confirm
$P(n)=e^{-|\alpha|^2}|\alpha|^{2n}/n!$ and $\bar n=(\Delta n)^2=|\alpha|^2$.
\begin{solution}
$P(n)=|\langle n|\alpha\rangle|^2=e^{-|\alpha|^2}|\alpha|^{2n}/n!$, the
Poisson distribution with parameter $|\alpha|^2$. Its mean and
variance both equal $|\alpha|^2$, so $\bar n=(\Delta n)^2=|\alpha|^2$
and $Q=0$. (Consistent with the ch.2 coherent-state exercise.)
\end{solution}
\end{exercise}

\begin{exercise}[Non-orthogonality]
Show $|\langle\beta|\alpha\rangle|^2=e^{-|\alpha-\beta|^2}$ and
comment on when two coherent states are nearly distinguishable.
\begin{solution}
Using the Fock expansion,
$\langle\beta|\alpha\rangle=e^{-|\alpha|^2/2-|\beta|^2/2}\sum_n(\beta^*\alpha)^n/n!=e^{-|\alpha|^2/2-|\beta|^2/2+\beta^*\alpha}$.
Then $|\langle\beta|\alpha\rangle|^2=e^{-|\alpha|^2-|\beta|^2+2\mathrm{Re}(\beta^*\alpha)}=e^{-|\alpha-\beta|^2}$.
They become nearly orthogonal (distinguishable) only when
$|\alpha-\beta|\gg1$, i.e.\ their phase-space separation exceeds the
vacuum noise radius.
\end{solution}
\end{exercise}

\begin{exercise}[Coherent state is minimum uncertainty]
Show $\Delta X=\Delta P=\tfrac12$ for any $\ket\alpha$.
\begin{solution}
$\langle\hat X\rangle=\mathrm{Re}\,\alpha$,
$\langle\hat X^2\rangle=\tfrac14\langle(\hat a+\hat a^\dagger)^2\rangle
=\tfrac14(\alpha^2+\alpha^{*2}+2|\alpha|^2+1)$ using
$\hat a\ket\alpha=\alpha\ket\alpha$; subtracting
$\langle\hat X\rangle^2=\tfrac14(\alpha+\alpha^*)^2$ leaves
$(\Delta X)^2=\tfrac14$. Identically $(\Delta P)^2=\tfrac14$. So the
displaced vacuum carries the same minimum, symmetric noise as the
vacuum.
\end{solution}
\end{exercise}

\begin{exercise}[Squeezing in decibels]
A squeezed vacuum reduces one quadrature variance by $10\,$dB.
Find $r$ and the photon number $\bar n=\sinh^2 r$.
\begin{solution}
$10\log_{10}e^{2r}=10$ dB means $e^{2r}=10$, so
$r=\tfrac12\ln10\approx1.15$. Then
$\sinh r=\sinh(1.15)\approx1.43$ and $\bar n=\sinh^2 r\approx2.0$.
A 10-dB squeezed vacuum thus contains about two photons on average
--- a substantial, technically demanding amount of squeezing.
\end{solution}
\end{exercise}

\begin{exercise}[Squeezed quadrature variances]
Derive \cref{eq:squeeze-variances} for $\theta=0$ from the
Bogoliubov relation \cref{eq:squeeze-bogoliubov}.
\begin{solution}
With $\theta=0$, $\op S^\dagger\hat a\op S=\hat a\cosh r-\hat a^\dagger\sinh r$,
so $\op S^\dagger\hat X\op S=\tfrac12\op S^\dagger(\hat a+\hat a^\dagger)\op S
=\tfrac12[(\hat a+\hat a^\dagger)\cosh r-(\hat a^\dagger+\hat a)\sinh r]
=\hat X(\cosh r-\sinh r)=\hat X e^{-r}$. Hence in the squeezed vacuum
$(\Delta X)^2=\langle0|\op S^\dagger\hat X^2\op S|0\rangle=e^{-2r}\cdot\tfrac14$.
Similarly $\op S^\dagger\hat P\op S=\hat P e^{+r}$ gives
$(\Delta P)^2=\tfrac14 e^{+2r}$.
\end{solution}
\end{exercise}

\begin{exercise}[Beam splitter preserves commutators]
Verify that the output modes \cref{eq:beam-splitter} satisfy
$\comm{\hat a_{\rm out}}{\hat a_{\rm out}^\dagger}=1$ and
$\comm{\hat a_{\rm out}}{\hat b_{\rm out}^\dagger}=0$.
\begin{solution}
$\hat a_{\rm out}=\hat a\cos\vartheta+\hat b\sin\vartheta$, so
$\comm{\hat a_{\rm out}}{\hat a_{\rm out}^\dagger}
=\cos^2\vartheta\comm{\hat a}{\hat a^\dagger}+\sin^2\vartheta\comm{\hat b}{\hat b^\dagger}
=\cos^2\vartheta+\sin^2\vartheta=1$.
$\hat b_{\rm out}=-\hat a\sin\vartheta+\hat b\cos\vartheta$, so
$\comm{\hat a_{\rm out}}{\hat b_{\rm out}^\dagger}
=-\cos\vartheta\sin\vartheta\comm{\hat a}{\hat a^\dagger}+\sin\vartheta\cos\vartheta\comm{\hat b}{\hat b^\dagger}=0$.
The transformation is canonical.
\end{solution}
\end{exercise}

\begin{exercise}[Hong--Ou--Mandel dip]
Send $\ket{1}_a\ket{1}_b$ into a 50:50 splitter
($\vartheta=\pi/4$). Show the output has no $\ket1_a\ket1_b$
component (the coincidence rate vanishes).
\begin{solution}
The input is $\hat a^\dagger\hat b^\dagger\ket{00}$. Under the
splitter $\hat a^\dagger\to(\hat a^\dagger+\hat b^\dagger)/\sqrt2$,
$\hat b^\dagger\to(-\hat a^\dagger+\hat b^\dagger)/\sqrt2$, so
$\hat a^\dagger\hat b^\dagger\to\tfrac12(\hat a^\dagger+\hat b^\dagger)(-\hat a^\dagger+\hat b^\dagger)
=\tfrac12(-\hat a^{\dagger2}+\hat b^{\dagger2})$ (the cross terms
$\hat a^\dagger\hat b^\dagger$ cancel). Acting on vacuum gives
$\tfrac1{\sqrt2}(-\ket{2}_a+\ket2_b)$: both photons exit together,
and the $\ket1_a\ket1_b$ coincidence amplitude is zero.
\end{solution}
\end{exercise}

\begin{exercise}[$g^{(2)}$ of coherent and thermal light]
Evaluate $g^{(2)}(0)$ for a coherent state and for thermal light
using \cref{eq:g2}.
\begin{solution}
Coherent: $\hat a\ket\alpha=\alpha\ket\alpha$ gives
$\langle\hat a^\dagger\hat a^\dagger\hat a\hat a\rangle=|\alpha|^4$
and $\langle\hat a^\dagger\hat a\rangle^2=|\alpha|^4$, so
$g^{(2)}(0)=1$. Thermal: using $Q=\bar n$ and $g^{(2)}=1+Q/\bar n$,
$g^{(2)}(0)=1+\bar n/\bar n=2$. Coherent light is uncorrelated;
thermal light bunches.
\end{solution}
\end{exercise}

\begin{exercise}[Antibunching of a single photon]
Show $g^{(2)}(0)=0$ for the Fock state $\ket1$ and explain why this
is impossible classically.
\begin{solution}
For $\ket1$, $\hat a\hat a\ket1=\hat a(\sqrt1\ket0)=0$, so
$\langle\hat a^\dagger\hat a^\dagger\hat a\hat a\rangle=0$ while
$\langle\hat a^\dagger\hat a\rangle^2=1$; thus $g^{(2)}(0)=0$.
Classically $g^{(2)}(0)=\langle I^2\rangle/\langle I\rangle^2\geq1$
by the Cauchy--Schwarz inequality for a non-negative intensity
$I$, so $g^{(2)}(0)<1$ has no classical-field description: detecting
the one photon empties the mode, and a second detector cannot fire
simultaneously.
\end{solution}
\end{exercise}
